% Parte 2: o Metodo 1, no mesmo estilo do operador.
% Fonte: nota do screened Poisson (02/10/2026, reports/investigating_v2/
% yan_screened_poisson_plan.pdf), Secs. 3-6.2 e 8; codigo
% core/screened.py, core/screened_front.py (ScreenedFront).
% Numeracao continua depois de (17).

\section{O Método 1}

\begin{frame}{O Método 1: fontes em $S$, screened Poisson, $\Delta$-Poisson}
  \small
  \renewcommand{\arraystretch}{2.5}
  \textbf{Definição}\\[-0.4em]
  \begin{tabular}{@{}l@{\hspace{0.6em}}r@{\;}l@{\hspace{1.4em}}l@{}}
    \nota{(18)} & $g(\boldsymbol{s})$ & $:= \mathrm{NN}\bigl(v(\boldsymbol{s});\boldsymbol{\theta}\bigr), \quad \boldsymbol{s}\in S$ & \nota{tempo de cada candidato em $S$}\\
    \nota{(19)} & $w(\boldsymbol{s})$ & $= \exp\!\left(-\dfrac{g(\boldsymbol{s}) - \min_S g}{\varepsilon}\right)$ & \nota{peso de cada fonte}\\
    \nota{(20)} & $u - \varepsilon^2\,\nabla\cdot(A\nabla u)$ & $= B_S\,w, \quad A = \U^{T}\U$ & \nota{screened Poisson}\\
    \nota{(21)} & $\q_g$ & $= -\dfrac{\nabla u}{\|\U\nabla u\|}$ & \nota{direção, no lugar de (11)}\\
    \nota{(22)} & $\Delta\varphi$ & $= \nabla\cdot\q_g$ & \nota{o operador, (8)--(10)}
  \end{tabular}
\end{frame}

\begin{frame}{O Método 1 no FEM}
  \small
  \renewcommand{\arraystretch}{2.5}
  \textbf{Discretização}\\[-0.4em]
  \begin{tabular}{@{}l@{\hspace{0.6em}}r@{\;}l@{\hspace{1.4em}}l@{}}
    \nota{(23)} & $H_\varepsilon$ & $= M + \varepsilon^2 K_A$ & \nota{$M$: lumped mass; $K_A$: rigidez com $A$}\\
    \nota{(24)} & $(\boldsymbol{b}_S)_s$ & $= a_s\, w_s, \quad s \in S$ & \nota{$a_s$: área de superfície do nó $s$}\\
    \nota{(25)} & $H_\varepsilon\,\boldsymbol{u}$ & $= \boldsymbol{b}_S$ & \nota{SPD, sem constante livre}\\
    \nota{(26)} & $\q_g$ & $=$ (21) em cada elemento, média nos nós & \nota{como $\p$ no Método 0}\\
    \nota{(27)} & $K\boldsymbol{\varphi}$ & $= \boldsymbol{b}(\q_g), \quad \varphi \leftarrow \varphi - \min\varphi$ & \nota{(9)--(10), sem mudança}
  \end{tabular}
\end{frame}

\begin{frame}{Por que $-\varepsilon\log u$ é um tempo}
  \small
  Longe de $S$, com $\psi = -\varepsilon\log u$, a equação (20) vira
  \[
    \nabla\psi^{T}\!A\,\nabla\psi \;-\; \varepsilon\,\nabla\cdot(A\nabla\psi) \;=\; 1
    \qquad\text{\nota{Eikonal viscosa; $\varepsilon \to 0$ dá (1)}}
  \]
  Com a função de Green do screened Poisson, $G_\varepsilon(\boldsymbol{x},\boldsymbol{s}) \approx C_\varepsilon\, e^{-T(\boldsymbol{s},\boldsymbol{x})/\varepsilon}$:
  \[
    -\varepsilon\log u(\boldsymbol{x}) \approx \operatorname*{softmin}_{\boldsymbol{s}\in S}\,
      \bigl[g(\boldsymbol{s}) + T(\boldsymbol{s},\boldsymbol{x})\bigr]
  \]
  \vspace{-0.6em}
  \begin{itemize}\setlength{\itemsep}{0.1em}\footnotesize
    \item cada candidato entra com o seu tempo $g(\boldsymbol{s})$, sem resolver a Eikonal;
    \item só tempos relativos contam: $g + C$ multiplica $u$ por $e^{-C/\varepsilon}$;
    \item as fontes ficam em $S$ por construção, sem termo $L_S$;
    \item $\nabla\times\q_g = \boldsymbol{0}$ não é garantido (normalização anisotrópica): o $\Delta$-Poisson continua necessário.
  \end{itemize}
  \vspace{-0.3em}
  {\scriptsize\color{gray} Nos resultados: $\varepsilon = 2$ ms, lumped mass, $H_\varepsilon$ fatorado uma vez com as velocidades fixas.}
\end{frame}

% Custo: problems/purkinje_roots/scripts/bench_mesh.py (.114, RTX 4070 Ti
% SUPER, 2026-10-08), via reports/investigating_v2/handoff_camps_papers.md.
\begin{frame}{Método 0 e Método 1}
  \small
  \renewcommand{\arraystretch}{1.35}
  \begin{tabular}{@{}p{0.25\textwidth}p{0.33\textwidth}p{0.36\textwidth}@{}}
    \toprule
    & Método 0 ($\Delta$-PoIssoNN) & Método 1 \\
    \midrule
    incógnita & $\p(\boldsymbol{x})$ em todo $\Omega$ (vetor) & $g(\boldsymbol{s})$ só em $S$ (escalar) \\
    fontes & emergem; precisam de um termo $L_S$ & só em $S$, por construção \\
    direção & $\p/\|\U\p\|$, qualquer campo & $-\nabla u/\|\U\nabla u\|$, de um escalar \\
    solves por iteração & 1 Poisson & 1 screened + 1 Poisson \\
    tempo por iteração (pac2) & 3,8 ms & 6,5 ms \\
    \bottomrule
  \end{tabular}
\end{frame}
